% ============================================================================
% DRAFT: e2_policy_family_table.tex   (W2 deliverable, E2 landing draft)
% Date: 2026-08-04 / package: 2026-08-05
% Source (frozen, read-only):
%   experiments/security-analysis-ablation-and-overhead/results/
%     policy-family-sensitivity/policy-family-sensitivity-report.md
%   Script: experiments/security-analysis-ablation-and-overhead/source/
%     policy-family-sensitivity/run_policy_family_sensitivity.py
%   Inputs: finite-contexts.jsonl (56 contexts, AgentDojo) +
%     heldout-contexts.jsonl (25 contexts, ToolSandbox)
% Status: INDEPENDENT DRAFT. Do NOT \input into main.tex until the PM
%         approves placement and the claim-map diff review has run.
%
% ----------------------------------------------------------------------------
% 落位建议 (placement suggestion)
%   File    : sections/results.tex
%   Position: new \subsection{Policy-Family Sensitivity} (D1 skeleton R.8),
%             AFTER \subsection{Pre-Registered Held-Out Contract Check} and
%             BEFORE \subsection{Common-Checkpoint AgentDojo Comparison}.
%   Note    : the three accompanying text paragraphs are already drafted in
%             defensive_writing_drafts_2026-08-04.md §C.2 (draft 2 gives the
%             claim-boundary paragraph with the frozen 7/10 and 20/24
%             summaries); land them together with this table.
%   Cross-refs: t1_theorem_section.tex and t3_separation_propositions.tex
%             both cite \ref{tab:policy-family-sensitivity}.
%
% ----------------------------------------------------------------------------
% Number policy (数字策略)
%   ALL numbers below are from the frozen E2 grid (48 deletion rows =
%   4 families x 6 qualifier settings x 2 domains, plus ceilings and
%   baseline-reproduction checks). They are permitted in LaTeX per the W2
%   constraint. No v17 / V0-V3 numbers appear in this file.
% ============================================================================

\begin{table*}[t]
\centering
\footnotesize
\setlength{\tabcolsep}{4pt}
\caption{Policy-family $\times$ qualifier-deletion sensitivity over two frozen
finite domains. For each domain--family--qualifier cell: \emph{Separating}
counts context pairs whose post-deletion representations coincide while some
authority in the family still separates their source effects
(authorization-separating collisions retained by the deletion);
\emph{Redundant} counts same-representation pairs whose full source effects
differ but whose family projection is identical (the distinction is
policy-invisible under that family); \emph{Overpart.} counts
family-equivalence classes split into more than one representation signature;
\emph{False rej.} counts family-equivalent pairs with different
representation signatures. Rows labeled \texttt{none} give each family's
no-deletion baseline. \textbf{Family enumeration principle.} The four
authority families are not selected from the data: they are the images of the
paper's submultiset authority construction under a pre-enumerated lattice of
occurrence projections---identity (the power-set baseline used throughout the
paper), the effect-name quotient, the resource quotient, and the
$\{0,1,\geq2\}$ count truncation of the effect-name quotient. Every family
contains every submultiset of every observed projected effect multiset, so
each family is closed under submultisets by construction. The family set was
fixed before any metric was computed and was not adjusted after observing
redundancy rates. \textbf{Both directions reported.} For every cell we report
both retained separating pairs (the qualifier is still doing work) and
redundant pairs (the qualifier is policy-invisible under that family),
together with overpartition and false-rejection counts; necessity claims are
scoped to cells marked \emph{necessary}, cells marked \emph{redundant}
delimit where the coarser family cannot express the distinction, and
\emph{vacuous} marks qualifiers not instantiated in that domain.
\textbf{Claim boundary (fixed scope).} Deterministic recomputation over two
frozen finite domains with source-hash-bound effect oracles; separability is
defined relative to the declared family's projection of source-effect
multisets. Results do not extend to unenumerated calls, open tool domains,
deployment policy languages, or families outside the enumerated lattice.}
\label{tab:policy-family-sensitivity}
\begin{tabular}{llrrrrl}
\toprule
Family & Qualifier deleted & Sep. & Red. & Overpart. & False rej. & Verdict \\
\midrule
\multicolumn{7}{l}{\textbf{(a) AgentDojo finite domain (56 contexts)}} \\
power\_set        & none       & 0  & 0  & 0 & 0   & no effect \\
                  & date       & 16 & 0  & 0 & 0   & necessary \\
                  & subject    & 16 & 0  & 0 & 0   & necessary \\
                  & recurrence & 16 & 0  & 0 & 0   & necessary \\
                  & payload    & 8  & 0  & 0 & 0   & necessary \\
                  & visibility & 4  & 0  & 0 & 0   & necessary \\
\midrule
effect\_kind      & none       & 0  & 0  & 6 & 548 & no effect \\
                  & date       & 0  & 16 & 6 & 532 & redundant \\
                  & subject    & 0  & 16 & 6 & 532 & redundant \\
                  & recurrence & 0  & 16 & 6 & 532 & redundant \\
                  & payload    & 0  & 8  & 5 & 540 & redundant \\
                  & visibility & 0  & 4  & 6 & 544 & redundant \\
\midrule
resource          & none       & 0  & 0  & 9 & 262 & no effect \\
                  & date       & 0  & 16 & 9 & 246 & redundant \\
                  & subject    & 0  & 16 & 9 & 246 & redundant \\
                  & recurrence & 0  & 16 & 9 & 246 & redundant \\
                  & payload    & 2  & 6  & 8 & 256 & necessary \\
                  & visibility & 0  & 4  & 9 & 258 & redundant \\
\midrule
count\_truncated  & none       & 0  & 0  & 6 & 548 & no effect \\
                  & date       & 0  & 16 & 6 & 532 & redundant \\
                  & subject    & 0  & 16 & 6 & 532 & redundant \\
                  & recurrence & 0  & 16 & 6 & 532 & redundant \\
                  & payload    & 0  & 8  & 5 & 540 & redundant \\
                  & visibility & 0  & 4  & 6 & 544 & redundant \\
\midrule
\multicolumn{7}{l}{\textbf{(b) ToolSandbox held-out domain (25 contexts)}} \\
power\_set        & none       & 0  & 0  & 0 & 0   & no effect \\
                  & date       & 4  & 0  & 0 & 0   & necessary \\
                  & subject    & 0  & 0  & 0 & 0   & vacuous \\
                  & recurrence & 0  & 0  & 0 & 0   & vacuous \\
                  & payload    & 16 & 0  & 0 & 0   & necessary \\
                  & visibility & 0  & 0  & 0 & 0   & no effect \\
\midrule
effect\_kind      & none       & 0  & 0  & 5 & 61  & no effect \\
                  & date       & 0  & 4  & 5 & 57  & redundant \\
                  & subject    & 0  & 0  & 5 & 61  & vacuous \\
                  & recurrence & 0  & 0  & 5 & 61  & vacuous \\
                  & payload    & 0  & 16 & 5 & 45  & redundant \\
                  & visibility & 0  & 0  & 5 & 61  & no effect \\
\midrule
resource          & none       & 0  & 0  & 3 & 57  & no effect \\
                  & date       & 0  & 4  & 3 & 53  & redundant \\
                  & subject    & 0  & 0  & 3 & 57  & vacuous \\
                  & recurrence & 0  & 0  & 3 & 57  & vacuous \\
                  & payload    & 0  & 16 & 3 & 41  & redundant \\
                  & visibility & 0  & 0  & 3 & 57  & no effect \\
\midrule
count\_truncated  & none       & 0  & 0  & 5 & 63  & no effect \\
                  & date       & 0  & 4  & 5 & 59  & redundant \\
                  & subject    & 0  & 0  & 5 & 63  & vacuous \\
                  & recurrence & 0  & 0  & 5 & 63  & vacuous \\
                  & payload    & 0  & 16 & 5 & 47  & redundant \\
                  & visibility & 0  & 0  & 5 & 63  & no effect \\
\bottomrule
\end{tabular}

\par\vspace{2mm}
\begin{tabular}{llrrr}
\toprule
\multicolumn{5}{l}{\textbf{(c) Family separability ceilings}} \\
Domain & Family & Diff.-effect pairs & Family-separable & Ceiling \\
\midrule
AgentDojo (56)      & power\_set       & 1540 & 1540 & 1.000 \\
                    & effect\_kind     & 1540 & 992  & 0.644 \\
                    & resource         & 1540 & 1278 & 0.830 \\
                    & count\_truncated & 1540 & 992  & 0.644 \\
\midrule
ToolSandbox (25)    & power\_set       & 468  & 468  & 1.000 \\
                    & effect\_kind     & 468  & 407  & 0.870 \\
                    & resource         & 468  & 411  & 0.878 \\
                    & count\_truncated & 468  & 405  & 0.865 \\
\bottomrule
\end{tabular}
\end{table*}

% ============================================================================
% NOTES FOR PORTING
%
% 1. Grid completeness: panels (a)+(b) contain all 48 deletion rows of the
%    frozen E2 report (4 families x {none,date,subject,recurrence,payload,
%    visibility} x 2 domains); panel (c) reproduces the separability ceilings.
% 2. Redundancy rate (not shown as a column to save width) is computable as
%    Red/(Red+Sep); the single partial exception is resource x payload on the
%    AgentDojo domain (Sep 2, Red 6, rate 0.75), visible directly in the row.
% 3. Baseline reproduction checks (from the frozen report; cite in Results
%    text, not in the table): typed contract under power-set/no-deletion
%    reproduces the zero-collision baseline on both domains; dropping all
%    qualifiers reproduces the frozen common-contract baselines of 118
%    (AgentDojo) and 41 (ToolSandbox) separating pairs; the vacuous
%    subject/recurrence deletions leave all ToolSandbox metrics unchanged.
% 4. Observed direction: MIXED (per-role verdicts: date/subject/recurrence
%    necessary only under power-set on AgentDojo; payload mixed; visibility
%    necessary only under power-set on AgentDojo; on ToolSandbox, date and
%    payload necessary only under power-set, subject/recurrence vacuous,
%    visibility never necessary under the tested families). The Results text
%    must use the mixed claim-boundary template; both alternative templates
%    (low-redundancy / high-redundancy) remain archived in the E2 report.
% 5. Suggested Results text paragraphs (family enumeration before results;
%    full grid both directions; mixed claim boundary) are drafted in
%    defensive_writing_drafts_2026-08-04.md §C.2 and use only frozen numbers
%    (7/10 power-set deletion rows create collisions; 20/24 non-vacuous
%    coarser-family rows have redundancy rate 1.000).
% ============================================================================
